\section{Attribution and proof boundaries}
The double-coset argument is a general normalization of a group action. The marked-inversion theorem makes explicit the chamber permutation behind the criterion in Fernandes--Leemans--Weiss (2016). The invariant kernel and mirror mix are standard constructions in the theory of chiral covers; Cunningham's 2023 author manuscript uses them. The intersection-transport proof has also been used in the earlier proof bank. Their value here is that the applications can invoke these fully stated interfaces without repeating the orbit and kernel arguments. No claim of historical priority is attached to this organization.

The retained geometric inputs are the string $C^+$ correspondence, its recursive intersection criterion, the marked-cover interpretation and the Schulte--Weiss theorem that facets of facets and vertex-figures of vertex-figures of a chiral polytope are directly regular. The last theorem is used in the application, not in the elementary kernel proof. The standard affine description of square torus maps and the specified older classification of atomic tight polytopes are additional application inputs, stated there.

Formal certificates cover only their declared algebraic and action statements. The seven new declarations in \texttt{TwoOrbitReflection.lean} check the actual two-copy formulas, including involutivity from the generating and inversion hypotheses. Two unchanged modules provide seven further compiled declarations; their declaration-to-argument map distinguishes the used intersection and propagation interfaces from unused statements in those files. They do not formalize the whole incidence-geometry definition, the string $C^+$ correspondence, the older classification or the six finite groups. The complete required GAP program for those groups is supplied separately. Its trusted exact group algorithms form a computational premise for that one finite example; the general theorems and the infinite-family deductions are written proofs. Later results on the existence or multiplicity of regular covers are separate results, not premises silently imported from the three earlier cover manuscripts.

\begin{thebibliography}{9}
\bibitem{FLW} M. E. Fernandes, D. Leemans and A. Ivi\'c Weiss, \emph{Highly symmetric hypertopes}, Aequationes Mathematicae 90 (2016), 1045--1067. Theorem 8.2.
\bibitem{ZL} W.-J. Zhang and D. Leemans, \emph{Testing chirality on hypertopes}, arXiv:2511.14562v1 (November 2025); subsequently Journal of Algebraic Combinatorics 63 (2026), article 61.
\bibitem{C} G. Cunningham, \emph{Chiral polytopes whose smallest regular cover is a polytope}, author manuscript, 25 July 2023; Journal of Combinatorial Theory, Series A 204 (2024), 105839. The checked edition is the eight-page author manuscript.
\bibitem{SW} E. Schulte and A. Ivi\'c Weiss, \emph{Chiral polytopes}, Applied Geometry and Discrete Mathematics, DIMACS Series 4 (1991), 493--516. The correspondence and section theorem are retained as cited by the source.
\bibitem{CP} G. Cunningham and D. Pellicer, \emph{Tight chiral polytopes}, Journal of Algebraic Combinatorics 54 (2021), 837--878. The specified Table 4 families and published Proposition 11 are retained; the acquired arXiv excerpt labels the latter Proposition 4.2.
\end{thebibliography}
